%!TEX root = ../hutbthesis_main.tex
\keywordscn{自动牵引车；飞机对接；ROS；Gazebo；路径规划；联合仿真}
\begin{abstractzh}

随着机场地面运行智能化水平的提升，飞机牵引与对接作业对自动化、精确性和安全性提出了更高要求。针对传统人工牵引方式依赖经验、作业效率受限及安全风险较高等问题，本文以自动牵引车与飞机对接过程为研究对象，基于ROS与Gazebo平台构建联合仿真系统。

本文首先对机场牵引作业场景进行抽象，建立双桥转向牵引车运动学模型，并分析转向角、速度与轨迹曲率之间的约束关系。在此基础上，设计飞机辅助感知与相对位姿误差控制方法，通过目标身份确认、相对距离与姿态误差获取，实现牵引车的自动接近、姿态调整和停车判定。随后，基于Gazebo搭建牵引车、飞机目标和机场作业环境，并通过ROS节点完成车辆控制、传感器数据获取、任务模式切换和系统通信。

仿真结果表明，所构建的ROS-Gazebo联合仿真系统能够实现牵引车模型加载、转向控制、轮速控制、目标接近、对接、推出和撤离等基本流程，验证了所提方法在机场自动牵引对接场景中的可行性。本文研究可为后续真实样机开发、控制算法优化和多传感器融合研究提供一定参考。

\end{abstractzh}
